\section{Experiments}\label{sec:experiments}

\emph{Setup.} All experiments were run on a server with four AMD EPYC 7513 32-core processors at 2.6\,GHz
and 504\,GB of memory.
Each instance was allocated a wall-clock timeout of 180\,seconds and a memory limit of 40\,GB.
We evaluated our solver, \solver, against Z3~4.15.4~\cite{z3}, cvc5~1.3.2~\cite{cvc5}, MathSAT5~5.6.16~\cite{mathsat5},
and Yices~2.6.4~\cite{Dutertre:cav2014} on the QF\_NIA benchmark set from SMT-LIB~\cite{smtlib}.
\solver is implemented in C++ using the smt-switch abstraction layer~\cite{mann2021smt},
linked against the same Z3~4.15.4 build used as the standalone Z3 baseline,
so any difference in results is due to the solving strategy rather than the underlying LIA engine.

Table~\ref{tab:results} summarises the results.
We report three configurations of \solver: the base solver (\solver), the base solver with
the frontier strategy for tangent-plane lemmas (\solver+\textsc{Frontier}), and an ablation
with lazy congruence axioms enabled (\solver+\textsc{Congr.}); the first two run without
congruence.

\begin{table}[t]
	\centering
	\setlength{\tabcolsep}{0.8em}
	\renewcommand{\arraystretch}{1.2}
	\begin{tabular}{lrrr|rrr}
		\hline
		                          & \multicolumn{3}{c|}{\textbf{Total} (25\,444)} & \multicolumn{3}{c}{\textbf{MathProblems} (1\,100)}                                                                     \\
		\textbf{Solver}           & \textbf{Sat}                                  & \textbf{Unsat}                                     & \textbf{Solved} & \textbf{Sat} & \textbf{Unsat} & \textbf{Solved} \\
		\hline
		\solver+\textsc{Frontier} & 8\,377                                        & 5\,709                                             & 14\,086         & 579          & 7              & 586             \\
		\solver                   & 8\,331                                        & 5\,680                                             & 14\,011         & 578          & 7              & 585             \\
		\solver+\textsc{Congr.}   & 6\,896                                        & 5\,525                                             & 12\,421         & 580          & 7              & 587             \\
		\hline
		Z3~4.15.4                 & 12\,718                                       & 6\,487                                             & 19\,205         & 111          & 7              & 118             \\
		MathSAT5~5.6.16           & 11\,177                                       & 5\,200                                             & 16\,377         & 164          & 7              & 171             \\
		Yices~2.6.4               & 10\,672                                       & 5\,181                                             & 15\,853         & 112          & 7              & 119             \\
		cvc5~1.3.2                & 8\,550                                        & 4\,446                                             & 12\,996         & 150          & 7              & 157             \\
		\hline
	\end{tabular}
	\vspace{0.9em}
	\caption{Instances solved within 180\,s on the SMT-LIB QF\_NIA benchmark set.\label{tab:results}}
\end{table}


\begin{table}[t]
    \centering
    \resizebox{\linewidth}{!}{%
    \begin{tabular}{l | lll | llll}
        \hline
        \textbf{Family} & \textbf{\solver+\textsc{Fr.}} & \textbf{\solver} & \textbf{\solver+\textsc{Cg.}} & \textbf{Z3} & \textbf{MathSAT} & \textbf{Yices} & \textbf{cvc5} \\
        \hline
        ITS (17046) & 4203/3908 & 4197/3881 & 3016/3738 & 8154/4452 & 6777/3569 & 6446/3434 & 4957/2843 \\
        AProVE (2409) & 1418/579 & 1397/579 & 1262/561 & 1642/687 & 1608/557 & 1590/708 & 1405/601 \\
        SAT14 (1926) & 1568/61 & 1566/61 & 1536/62 & 1853/72 & 1722/66 & 1809/65 & 1573/72 \\
        CInteger (1818) & 387/580 & 373/576 & 294/575 & 704/655 & 652/435 & 503/448 & 262/368 \\
        ReachSafety-Loops (350) & 9/310 & 9/310 & 9/317 & 11/339 & 10/320 & 6/302 & 11/322 \\
        mcm (186) & \textbf{11/0} & \textbf{12/0} & \textbf{11/0} & 8/0 & 7/0 & 6/0 & 9/0 \\
        calypto (177) & \textbf{80/97} & \textbf{80/97} & \textbf{80/97} & 79/97 & 79/90 & 79/95 & 79/96 \\
        leipzig (167) & 96/1 & 94/1 & 81/1 & 131/1 & 127/2 & 100/1 & 78/2 \\
        LassoRanker (106) & 4/91 & 4/92 & 4/91 & 4/100 & 4/101 & 4/85 & 4/92 \\
        UltimateAutomizerSvcomp2023 (58) & \textbf{7/15} & \textbf{7/15} & \textbf{7/15} & 8/12 & 7/10 & 5/0 & 6/0 \\
        UltimateLassoRanker (32) & 6/26 & 6/26 & 6/26 & 6/26 & 6/26 & 6/26 & 6/26 \\
        sqrtmodinv-hoenicke (27) & 0/17 & 0/17 & 0/17 & 0/17 & 0/0 & 0/0 & 0/1 \\
        ConcurrencySafety-Main (24) & 5/8 & 5/9 & 5/9 & 2/14 & 7/9 & 0/3 & 4/8 \\
        elster (9) & 4/0 & 3/0 & 5/0 & 5/0 & 7/0 & 6/0 & 6/0 \\
        UltimateAutomizer (7) & 0/7 & 0/7 & 0/7 & 0/7 & 0/7 & 0/7 & 0/7 \\
        LCTES (2) & \textbf{0/2} & \textbf{0/2} & \textbf{0/2} & 0/1 & 0/1 & 0/0 & 0/1 \\
        \hline
    \end{tabular}%
    }
    \vspace{0.9em}
    \caption{Sat/unsat instances solved per benchmark family within 180\,s (excluding \texttt{20220315-MathProblems}, see Table~\ref{tab:results}).\label{tab:families}}
\end{table}


On the full benchmark set \solver is competitive with cvc5 and trails Z3, MathSAT, and
Yices. Z3 combines a large number of built-in tactics and heuristics; MathSAT is
closed source. \solver, by contrast, is a concise open-source implementation with a
single LIA backend and no problem-specific tuning.

Enabling the frontier strategy adds a small but consistent improvement at no cost.
Enabling congruence axioms (\solver+\textsc{Congr.}) reduces the overall count by around
1\,600 instances, confirming that the overhead of pairwise congruence checks outweighs
their benefit on the general benchmark set.

The gap on satisfiable instances is most visible in large families such as ITS, AProVE, and
SAT14 (Table~\ref{tab:families}), and is largely structural: \solver must witness satisfiability
by iterative axiom refinement over a LIA abstraction, and each candidate solution is validated
purely through linear arithmetic.
Competing solvers employ complementary techniques, such as bit-blasting to fixed-width
integers, that can more directly enumerate satisfying
assignments for large-variable instances without the overhead of abstraction-refinement.
These families are particularly amenable to bit-blasting: the nonlinear structure is shallow,
instances are large (routinely hundreds of pures), and satisfying assignments tend to involve
small values.

The \texttt{20220315-MathProblems} family consists of 1100 instances encoding number-theoretic
problems, many of which involve sums of cubes ($x^3 + y^3 + z^3 = n$).
\solver solves 585--587 of these (53\%). This gap directly reflects the contribution of the
secant-based axioms for $x^k$: the convergence on cubic constraints requires bounds that go
beyond tangent-plane lemmas.
